% Generated draft: requires mathematical and source review.
\documentclass[11pt]{article}
\usepackage[margin=25mm]{geometry}
\usepackage{fontspec}
\setmainfont{Latin Modern Roman}
\usepackage{amsmath,amssymb,mathtools,unicode-math}
\setmathfont{Latin Modern Math}
\usepackage{enumitem,needspace,xurl,hyperref}
\setlength{\parindent}{0pt}
\setlength{\parskip}{5pt}
\setlength{\emergencystretch}{3em}
\newcommand{\N}{\mathbb{N}}
\newcommand{\Z}{\mathbb{Z}}
\newcommand{\Q}{\mathbb{Q}}
\newcommand{\R}{\mathbb{R}}
\newcommand{\C}{\mathbb{C}}
\newcommand{\F}{\mathbb{F}}
\newcommand{\A}{\mathbb{A}}
\newcommand{\PP}{\mathbb{P}}
\newcommand{\E}{\mathbb{E}}
\newcommand{\eps}{\varepsilon}
\newcommand{\dd}{\,\mathrm{d}}
\begin{document}
\section*{UnsolvedMath ID: 700279}
\textbf{Source problem code:} catalogue entry
\subsection{Definitions and mathematical statement}
A variety of groups is a class defined by group identities. For a variety \(\mathcal V\) and a set \(X\), let \(F_{\mathcal V}(X)\) be the relatively free group in \(\mathcal V\) on \(X\). For a group \(G\), define its derived series by \(G^{(0)}=G\) and \(G^{(n+1)}=[G^{(n)},G^{(n)}]\). Call \(\mathcal V\) soluble if there exists an integer \(d\geq 0\) such that \(G^{(d)}=1\) for every \(G\in\mathcal V\).

A group \(G\) is residually torsion-free nilpotent if, for every \(g\in G\setminus\{1\}\), there exist a torsion-free nilpotent group \(N\) and a homomorphism \(\varphi:G\to N\) with \(\varphi(g)\neq 1\). Here torsion-free means that no nonidentity element has finite order, and nilpotent means that the lower central series \(\gamma_1(N)=N\), \(\gamma_{n+1}(N)=[\gamma_n(N),N]\) reaches the trivial subgroup after finitely many steps.

\textbf{Group-variety question (A. L. Shmel'kin, 1973).} Under the interpretation of the source phrase ``residually nilpotent torsion-free'' as residually torsion-free nilpotent, is the following implication true for every variety of groups \(\mathcal V\)?
\[
\left(\forall\text{ sets }X,\quad F_{\mathcal V}(X)\text{ is residually torsion-free nilpotent}\right)
\ \Longrightarrow\ 
\left(\mathcal V\text{ is soluble}\ \text{or}\ \mathcal V=\mathcal G\right),
\]
where \(\mathcal G\) is the variety of all groups. The source wording does not explicitly distinguish this residual property from being both torsion-free and residually nilpotent; these are different hypotheses. Under the latter reading, the premise instead requires each \(F_{\mathcal V}(X)\) to be torsion-free and to satisfy \(\bigcap_{n\geq 1}\gamma_n(F_{\mathcal V}(X))=1\). The intended reading requires source review.

\textbf{Supplied sufficient Lie-algebra assertion.} A variety of Lie algebras over \(\Q\) is a class defined by Lie-algebra identities over \(\Q\). For a Lie algebra \(L\), put \(L^{(0)}=L\) and \(L^{(n+1)}=[L^{(n)},L^{(n)}]\). Such a variety \(\mathcal W\) is soluble if there exists an integer \(d\geq 0\) such that \(L^{(d)}=0\) for every \(L\in\mathcal W\). With the convention that a simple Lie algebra is nonabelian and has no ideals other than \(0\) and itself, the auxiliary assertion is
\[
\begin{gathered}
\text{for every variety }\mathcal W\text{ of Lie algebras over }\Q,\\
\left(\nexists S\in\mathcal W:\ \dim_{\Q}S<\infty\text{ and }S\text{ is simple}\right)
\ \Longrightarrow\ \mathcal W\text{ is soluble}.
\end{gathered}
\]
The source states that proving this Lie-algebra assertion suffices for an affirmative answer to the group-variety question; it does not state an equivalence. The source does not explicitly specify its convention for simplicity, so the nonabelian convention used above requires confirmation. Neither assertion is proved here.
\subsection{Short English statement}
The question asks whether a variety whose relatively free groups have the stated residual property must have a uniform bound on derived length or be the variety of all groups. The source also gives a sufficient assertion about varieties of Lie algebras over \(\Q\) containing no finite-dimensional simple members. The residual-property wording and the convention for simple Lie algebras require review. The supplied registry labels the problem open; that status has not been independently verified.
\subsection{Sources}
\begin{itemize}
\item Local UnsolvedMath registry, supplied complete record: permanent ID \(700279\), exact source code \texttt{catalogue entry}, title \emph{Kourovka Notebook Problem 4.72}. The supplied status is \texttt{open}.
\item \emph{The Kourovka Notebook (Unsolved Problems in Group Theory)}, Problem 4.72, attributed to A. L. Shmel'kin, 4th issue (1973). The supplied background identifies the reference edition as the 21st issue edition (2026), \texttt{arXiv:1401.0300v48}. This bibliographic attribution is supplied source data and has not been independently checked.
\item The supplied collection description attributes editorship to E. I. Khukhro and V. D. Mazurov and describes coverage of Issues 1--21 (1965--2026). The record carries upstream revision \texttt{d371ee917ac9647fbb06cd6fc529deec6f3dda82} and an update timestamp of \texttt{2026-10-16T14:00:00Z}; this timestamp is registry metadata, not a literature-verification date.
\end{itemize}
\end{document}
